فرض کنید $S$ حاصل‌ضرب همه‌ی اعداد اول کوچک‌تر از $2014$ باشد. بنابر قضیه‌ی دیریکله، بی‌نهایت عدد اول $p$ وجود دارد به طوری که $p \equiv 1 \pmod{S}$.

با استفاده از قضیه‌ی باقیمانده‌ی چینی، عدد طبیعی $n$ وجود دارد که
\[
n \equiv -c \pmod{p}, \qquad n \equiv 1 \pmod{p-1}
\]
(و در صورت نیاز $n \equiv 1 \pmod{q}$ برای همه‌ی اعداد اول $q < 2014$ که $p-1$ را نمی‌شمارند). آنگاه $\gcd(n,S)=1$، پس $n$ بر هیچ عدد اول کوچک‌تر از $2014$ بخش‌پذیر نیست. بنابر قضیه‌ی کوچک فرما به دست می‌آوریم
\[
p \mid n+c \mid a^n + b^n + n \equiv a + b + n \pmod{p}.
\]
چون $n \equiv -c \pmod{p}$ این نتیجه می‌دهد $a + b - c \equiv 0 \pmod{p}$. چون $p$ را می‌توان به‌دلخواه بزرگ انتخاب کرد، نتیجه می‌گیریم
\[
a + b = c. \tag{$\star$}
\]

اکنون عدد اول دلخواه $q > 2015$ را ثابت بگیرید. دوباره بنابر قضیه‌ی باقیمانده‌ی چینی عدد طبیعی $t$ وجود دارد به طوری که
\[
t \equiv -c \pmod{p}, \qquad t \equiv q \pmod{p-1}
\]
(و $\gcd(t,S)=1$). آنگاه
\[
p \mid t+c \mid a^t + b^t + t \equiv a^q + b^q - c \pmod{p}.
\]
چون $p$ می‌تواند به‌دلخواه بزرگ باشد، به دست می‌آوریم
\[
a^q + b^q = c. \tag{$\star\star$}
\]

از $(\star)$ و $(\star\star)$ نتیجه می‌شود که $a = b = 1$ و $c = 2$.

(به‌راحتی بررسی می‌شود که سه‌تایی $(1,1,2)$ واقعاً در شرط اصلی صدق می‌کند.)
